\sourceimport{finals-dp-extra}{Prelude.h}
\section{Three reusable DP and tree patterns}
These small models cover different shapes of dependency: a subset of
completed tasks, a parent-child choice, and a contiguous subtree range.
First identify which previous decisions can affect the next one. If that
information is a subset, use subset DP; if choices only interact across
tree edges, root the tree; if every query asks about descendants, flatten
the tree before selecting a range structure.

\subsection{Order every item once: subset Hamiltonian path}
Suppose each ordered pair of distinct jobs has a transition cost, and all
jobs must be done once with the least total transition cost. A greedy
``take the cheapest next job'' can block a good continuation. For example,
with costs $0\to1=1$, $0\to2=2$, $1\to2=100$, $2\to1=1$,
and both edges back to 0 costing 100, greedy
from 0 chooses 1 and pays 101; the order $0,2,1$ costs 3. The missing
information is which jobs have already been used.

Define $dp[S][v]$ as the least cost of an ordering of exactly the set
$S$, ending at $v$. Its last transition comes from a path on
$S\setminus\{v\}$ ending at some $u$:
\[
 dp[S][v]=\min_{u\in S\setminus\{v\}}
       (dp[S\setminus\{v\}][u]+c[u][v]).
\]
For a singleton, $dp[\{v\}][v]=0$: the start is unrestricted and no
edge has yet been paid. The answer is $\min_v dp[All][v]$, since the
endpoint is also unrestricted. On the example, $dp[\{0,2\}][2]=2$,
then $dp[All][1]=3$. This is the Held--Karp subset recurrence: $O(n^2 2^n)$
time and $O(n2^n)$ memory. Store predecessor vertices if the order is
needed. The tested helper returns only the cost and enforces $1\le n\le18$.

\sourceimport{finals-dp-extra}{SubsetOrder.h}
This model changes if the route must start at a fixed vertex (initialize
only that singleton), must return to start (add the closing edge), or may
skip jobs (change the objective and permitted terminal masks; with nonnegative costs,
freely skipping every job makes the optimum trivial). Precedence constraints require
rejecting transitions that place a job before its prerequisites. For
larger $n$, the exponential state count is the limit; use problem structure
such as bitmask width, meet-in-the-middle, or a different optimization.

\subsection{Pair disjoint endpoints in a tree}
For a tree, maximize the number of edges chosen so no two share an endpoint.
This is maximum matching, a useful starter when the statement says
``pair people/items across links'' and each endpoint can be used once. A
global greedy choice is risky: a locally attractive edge can consume a
vertex needed by multiple later choices. Root the tree; choices in
different child subtrees become independent once we record whether the
root is already paired with one child.

Let $free[u]$ be the best matching in the subtree when $u$ is not matched
to a child; let $best[u]$ allow either status. Then
\[
 free[u]=\sum_{v child}best[v],\qquad
 best[u]=\max\left(free[u],\max_{v child}
   \left(1+free[v]+\sum_{w\ne v}best[w]\right)\right).
\]
If $u$ is paired with $v$, that edge contributes one and $v$ cannot be
paired below itself, which explains $free[v]$. For a root with two leaf
children, $free=0$ and $best=1$: choose either edge. Each child is
considered as the one partner, so the recurrence is linear in total tree
degree. Supply symmetric adjacency for a connected simple tree. The implementation checks
endpoints, arc count, cycles, and connectivity, and
builds parent/order iteratively, then computes states in reverse order;
this avoids recursion depth failures on a chain.

\sourceimport{finals-dp-extra}{TreeMatching.h}
This assumes every edge has equal value and the goal is cardinality. With
edge weights, replace the added one with that edge's weight. If vertices
have capacities above one, the two-state recurrence no longer captures
their remaining capacity. If the graph has cycles, child subproblems are
not separated by one parent edge; use a graph matching algorithm instead.

\subsection{Turn descendants into a range}
When many queries ask for the sum, count, minimum, or updates over a node's
subtree, answering each by walking descendants can cost $O(n)$ per query.
In a rooted tree, a preorder DFS visits each subtree consecutively. Record
the position $tin[u]$ when $u$ is appended, and $tout[u]$ as the exclusive
end after all descendants. Then subtree $u$ is exactly the half-open range
$[tin[u],tout[u])$ in the flattened order. For example, root 0 with
children 1 and 2, where 1 has child 3, may flatten to $[0,1,3,2]$:
subtree 1 is positions $[1,3)$ and subtree 2 is $[3,4)$.

\sourceimport{finals-dp-extra}{SubtreeRanges.h}
Use a Fenwick tree for point updates and subtree sums: add a node's value
at $tin[u]$, and query the flattened interval. Offline threshold counts
can sort both events and nodes by value; activate nodes whose value meets
the current threshold in a Fenwick tree over $tin$, then ask the subtree
range sum. For example, if flattened values are $[5,2,7,4]$, the query
``count values at least 5 in subtree 1'' activates positions 0 and 2 and
returns the sum on $[1,3)$, namely 1.

The flattening supports any range structure, but it does not make arbitrary
path queries contiguous. Fixed-root subtree queries and path queries are
different models. Re-rooting queries also change which vertices are descendants.

\subsection{Find the kth object by skipping completion blocks}
Synthetic task: list length-$n$ binary strings with no adjacent ones in
lexicographic order and return rank $k$, starting from rank 1. Leading zeroes
are allowed. Do not generate all strings: count how many completions begin
with 0, then decide whether to stay in that block or skip it.

Let $W[L]$ count valid strings of length $L$ with no preceding 1.
$W[0]=1$, $W[1]=2$, and $W[L]=W[L-1]+W[L-2]$: start with 0 and any
valid suffix, or start with 10 and a shorter suffix. During reconstruction,
a preceding 1 forces 0. Otherwise the 0-block has $W[L-1]$ completions.
If $k$ is at most that count, choose 0. Otherwise subtract the count and
choose 1. Each prefix defines a contiguous block in lexicographic order.

For $n=3$, the list is $000,001,010,100,101$. At rank 4, the 0-block has
three entries, so skip it and choose 1 with new rank 1. Then 0 is forced;
the final digit is 0. The answer is 100.
\sourceimport{finals-dp-extra}{Unranking.h}
Counts used for rank comparisons must not be reduced modulo a prime. The
helper caps every count at the original requested rank using overflow-safe
addition. A capped count still tells whether the current rank fits, because
the current rank never increases. Cap at the original rank, not at a changing
rank while building the table. The helper returns no answer for an invalid
rank and handles $n=0$ as one empty string.

For a different object, derive its completion-count state first. Ordering
permutations by inversion count and then lexicographically, as in 2024 I,
needs a different recurrence and proved bounds; this binary-string helper
is not its solution. The reusable step is counting, skipping, and updating
the residual rank. Test reconstruction by enumerating all small objects.
